\chapter{Scope, conventions, and theorem contracts}

This project formalizes Miquel Ortega and Sean Prendiville,
\emph{Extremal Sidon sets are Fourier uniform, with applications to
partition regularity}, arXiv:2110.13447v1 (2021),
\url{https://arxiv.org/abs/2110.13447}.
The source is distributed under CC BY 4.0,
\url{https://creativecommons.org/licenses/by/4.0/}.
These expanded proofs distinguish mathematical arguments from
kernel-checked declarations. Only nodes explicitly marked as formalized
have checked Lean proofs. Proposed names remain ordinary text until
their declarations exist.

\section{Conventions}
Throughout $N\geq1$ is an integer, $[N]=\{1,\ldots,N\}\subseteq\mathbb Z$,
and an integer interval of length $L\in\mathbb N$ means
$\{a+1,\ldots,a+L\}$. Set
\[
 m=|S|,\qquad D_N(S)=\left|m/\sqrt N-1\right|+N^{-1/6}.
\]
Real powers apply to positive real casts of $N$. The empty Sidon set is
allowed. Division by $m$ requires an explicit positive lower bound.
A coloring has $r\geq1$ colors and is a function on $S$; its color
classes are disjoint and cover $S$. Solutions allow repeated coordinates.

For finitely supported functions on $\mathbb Z$, use counting measure and
\[
 e(t)=\exp(2\pi it),\quad
 \widehat f(\alpha)=\sum_x f(x)e(\alpha x),\quad
 (f*g)(x)=\sum_y f(y)g(x-y).
\]
There are no hidden normalization factors. The torus
$\mathbb T=\mathbb R/\mathbb Z$ carries probability Haar measure.
In Lean use \texttt{Finsupp} for finite-support functions and
\texttt{AddCircle 1} for the torus. Prove combinatorial
identities with finite sums before adding analytic interfaces.
An absolute implicit constant must be uniform in dimension $d$;
$O_d$ explicitly allows dimension dependence.
Implement asymptotic estimates with quantified constants and thresholds.

\section{Public targets}
\begin{theorem}[Paper Theorem 1.2]\label{target:uniformity}
\lean{Sidon.sidon_fourier_uniformity}\leanok
\uses{core:fourier-uniformity}
There exists $C>0$ such that for every $N\geq1$, every Sidon
$S\subseteq[N]$, and every $\alpha\in\mathbb T$,
\[
 \left|\widehat{1_S}(\alpha)-\frac mN\widehat{1_{[N]}}(\alpha)\right|
 \leq C\sqrt N D_N(S)^{1/2}.
\]
\end{theorem}
The pointwise estimate is equivalent to its supremum-norm version:
the uniform upper bound bounds the range, and every range element is
at most its supremum. No unproved boundedness hypothesis is suppressed.

\begin{corollary}[Largest sets, paper Corollary 1.4]\label{target:largest}
\lean{Sidon.largest_sidon_fourier_uniformity}\leanok
For sufficiently large $N$, a maximum-cardinality Sidon subset of $[N]$
satisfies
\[
 \left\|\widehat{1_S}-\frac mN\widehat{1_{[N]}}\right\|_\infty
 \ll \|\widehat{1_S}\|_\infty N^{-1/12}.
\]
\end{corollary}
\begin{proof}
A maximum exists in the finite power set of $[N]$, since the empty set
is admissible. The proved prime-interval theorem and Bose--Chowla
construction give $m\geq\sqrt N-128N^{1/3}$ for all sufficiently large
$N$. Appendix A gives $m\leq\sqrt N+4N^{1/4}$.
Dividing by $\sqrt N$ and using $N\geq1$ yields
\[
 |m/\sqrt N-1|\leq132N^{-1/6}.
\]
Theorem~\ref{target:uniformity} therefore gives the bound
$16\sqrt{133}\sqrt N\,N^{-1/12}$.
The lower size bound also gives $m\geq\sqrt N/2$ eventually.
The triangle inequality gives $|\widehat{1_S}(\alpha)|\leq m$,
and equality holds at zero, so $\|\widehat{1_S}\|_\infty=m$.
Substitute $\sqrt N\leq2m$ to obtain the absolute constant
$32\sqrt{133}$.
The full prime-interval proof and threshold choices are recorded in
Theorem~\ref{construction:largest-unconditional}.
\end{proof}

\begin{theorem}[Paper Theorem 1.15]\label{target:partition}
\lean{Sidon.sidon_partition_regularity}
\leanok
\uses{app:partition-regularity}
Fix $s\geq5$, $r\geq1$, and nonzero integers $c_1,\ldots,c_s$.
Suppose some nonempty subset of the coefficients sums to zero.
There exist $N_0\in\mathbb N$, $\eta>0$, and $\delta>0$, depending
only on $c$ and $r$, such that for every $N\geq\max(1,N_0)$ and Sidon
$S\subseteq[N]$ with $\big|m-\sqrt N\big|\leq\eta\sqrt N$,
every $r$-coloring has a class $C_j$ satisfying
\[
 \#\{x\in C_j^s:\textstyle\sum_i c_ix_i=0\}
 \geq\delta m^s/N.
\]
\end{theorem}
Shrink $\eta$ below $1/2$ to ensure $m>0$. The right side is positive,
so the counting conclusion implies existence. This is the precise
constant-and-threshold form of the paper's three alternatives.

\section{Implementation structure}
The main chain runs from Sidon differences through finite Fourier analysis,
interval bounds, equidistribution, dense models, and counting.
The improved Fourier branch depends on Cilleruelo's estimate.
The \href{https://github.com/TomasOrtega/Sidon/blob/master/docs/core-formalization-plan.md}{core plan}
and \href{https://github.com/TomasOrtega/Sidon/blob/master/docs/applications-formalization-plan.md}{applications plan}
map statements to exact declarations and files. The module structure is
\begin{verbatim}
Sidon/Basic.lean
Sidon/{Combinatorics,Analysis,Equidistribution,Transference}/
Sidon/Analysis/Approximation/
Sidon/{Constructions,NumberTheory}/
Sidon/{FourierUniformity,ImprovedFourierUniformity}.lean
Sidon/PartitionRegularity.lean
\end{verbatim}
